\documentclass[11pt]{article}
\usepackage[T1]{fontenc}
\usepackage{lmodern}
\usepackage{amsmath,amssymb,amsthm,mathtools}
\usepackage[left=1in,right=1in,top=0.85in,bottom=0.85in]{geometry}
\usepackage{microtype,booktabs,array}
\usepackage{xurl}
\usepackage[hidelinks,pdftitle={Unimodality of nested-box differences and near-full q-rook polynomials},pdfauthor={Aurantium},pdfsubject={A proof via nested-box differences}]{hyperref}
\usepackage{fancyhdr}
\pagestyle{fancy}
\fancyhf{}
\fancyhead[L]{Near-full $q$-rook unimodality}
\fancyhead[R]{Aurantium \textbar\ 13 October 2026}
\fancyfoot[C]{\thepage}
\setlength{\headheight}{14pt}
\setlength{\parskip}{0.35em}
\setlength{\emergencystretch}{2em}
\newtheorem{theorem}{Theorem}
\newtheorem{lemma}{Lemma}
\newtheorem{proposition}{Proposition}
\newtheorem{corollary}{Corollary}
\newcommand{\NN}{\mathbb Z_{\ge0}}
\newcommand{\inv}{\operatorname{inv}}
\newcommand{\supp}{\operatorname{supp}}
\newcommand{\qint}[1]{[#1]_q}
\newcommand{\BoxP}{P}
\title{Unimodality of nested-box differences\\
and near-full $q$-rook polynomials}
\author{Aurantium}
\date{13 October 2026}
\begin{document}
\maketitle
\thispagestyle{fancy}

\section*{Introduction}
For every integer $\ell\ge1$ and partition
$\lambda=(\lambda_1\ge\cdots\ge\lambda_\ell>0)$, let
\[
 B_\lambda=\{(i,j):1\le i\le\ell,\ 1\le j\le\lambda_i\}.
\]
A placement is a set of cells with at most one cell in each row and column.
A rook $(a,b)$ cancels the cells $(i,b)$ with $i\le a$ and $(a,j)$ with
$j\le b$. For a placement $w$, let $\inv(w)$ count the cells not cancelled
by any rook, and set
\[
 R_r(\lambda;q)=\sum_{\substack{w\subseteq B_\lambda\ \text{a placement}\\ |w|=r}}
 q^{\inv(w)}\in\NN[q].
\]
Here $q$ is an indeterminate. A coefficient sequence is unimodal when it
weakly increases and then weakly decreases; zero polynomials and zero
padding are allowed. The target is
\[
 \boxed{R_{\ell-1}(\lambda;q)\text{ is unimodal for every such }\lambda.}
\]
This is the final question in Lewis--Morales, arXiv:2410.19714v3,
Remark 6 (page 5); their rook conventions are on pages 1--2.
The setting is finite integer combinatorics. No field, characteristic,
regularity, or asymptotic hypothesis is required.

\section*{Main result}
The near-full polynomial is unimodal for every Ferrers board. The proof
establishes a stronger statement about nested products of finite chains.
It uses neither a Lefschetz property nor log-concavity of the rook polynomial,
and does not depend on finite computations. The product-difference reduction
and Hilbert-series interpretation originated in unpublished preliminary
research notes; both are derived explicitly below.

\begin{theorem}[Nested-box difference]\label{thm:boxes}
Let $n\ge1$ and $u,v\in\NN^n$ satisfy $v_i\le u_i$ for every $i$. Put
\[
 P_u(q)=\prod_{i=1}^n(1+q+\cdots+q^{u_i}).
\]
Then $P_u(q)-P_v(q)$ has nonnegative, unimodal coefficients.
\end{theorem}
The theorem includes $u=v$, where the difference is zero. Its application to
near-full rook polynomials has $u_i=a_i$ and $v_i=a_i-1$, after separating
negative and zero $a_i$. The proof is entirely elementary; all auxiliary
inequalities and the chain decomposition used below are proved here.

\subsection*{1. Checking the rook reduction and boundary cases}
Write the row lengths increasingly as $b_1\le\cdots\le b_\ell$ and set
$a_i=b_i-i+1$. Write $\qint m=1+q+\cdots+q^{m-1}$ for $m\ge1$ and
$\qint0=0$.

\begin{proposition}[Near-full product difference]\label{prop:rook}
If some $a_i<0$, then $R_{\ell-1}(\lambda;q)=0$. Otherwise,
\begin{equation}\label{eq:rook}
 R_{\ell-1}(\lambda;q)=\prod_{i=1}^{\ell}\qint{a_i+1}
                         -\prod_{i=1}^{\ell}\qint{a_i}.
\end{equation}
\end{proposition}
\begin{proof}
If $a_i<0$, the $i$ shortest rows have at most $b_i\le i-2$ columns
available. A placement with only one omitted row must put at least $i-1$
rooks into those rows, which is impossible.

Assume henceforth that every $a_i\ge0$. Process rows from the shortest
(bottom) row to the longest (top) row, and fix the unique omitted row $i$.
Every column used earlier exists in the current row and has already been
cancelled there. In a row with $w$ remaining columns, placing a rook leaves
$0,1,\ldots,w-1$ uncancelled cells, once each as the rook position varies.
Its weight sum is $\qint w$. No later rook can cancel a cell in this lower
row. Thus choices may be summed row by row, since the number $w$ is
independent of the previous column choices.

Before row $i$, the available widths are $a_j$; after row $i$, they are
$a_j+1$. The omitted row itself contributes $a_i$ uncancelled cells. Its
contribution is therefore
\[
 q^{a_i}\prod_{j<i}\qint{a_j}\prod_{j>i}\qint{a_j+1}.
\]
A width zero gives a zero factor, correctly excluding that placement type.
Summing over $i$, and using $\qint{a_i+1}-\qint{a_i}=q^{a_i}$, gives the
telescoping product identity \eqref{eq:rook}.
\end{proof}
If some $a_i=0$, the second product vanishes and the first is a product of
$q$-integers. Such products are symmetric and unimodal, as also follows
from the lemmas below. If all $a_i\ge1$, Theorem~\ref{thm:boxes} applies
to $u=a$ and $v=a-(1,\ldots,1)$. This covers every board and every $\ell$.
In particular, $\ell=1$ gives $R_0(\lambda;q)=q^{\lambda_1}$.

The unrestricted positive-vector problem is not an artificial enlargement:
any multiset of positive $a_i$ can be sorted increasingly and realized by
$b_i=a_i+i-1$, then reversing the $b_i$ to make a partition.
The basic non-log-concave example remains compatible with the result:
$a=(1,3)$, corresponding to $\lambda=(4,1)$, gives
\[
 R_1((4,1);q)=q+q^2+2q^3+q^4.
\]
Its unimodality is not log-concavity, since $1^2<1\cdot2$ at degree 2.

\subsection*{2. Two elementary interval-sum inequalities}
For a finitely supported sequence $h=(h_i)_{i\in\mathbb Z}$ of nonnegative
real numbers, log-concavity means $h_i^2\ge h_{i-1}h_{i+1}$ for all $i$.
``No internal zeros'' means its positive support is an integer interval.

\begin{lemma}[Sliding and shrinking sums]\label{lem:interval}
Suppose $h$ is log-concave with no internal zeros. For any integer $S$ and
integer $m\ge0$, both sequences
\[
 U_k=\sum_{i=k-m}^{k}h_i,
 \qquad
 V_k=\sum_{i=k-m}^{S-k}h_i
\]
are log-concave, with an empty sum interpreted as zero.
\end{lemma}
\begin{proof}
Fix an integer $t$. The sequence $w_t(i)=h_i h_{t-i}$ is symmetric about
$t/2$ and nondecreasing as $i$ approaches $t/2$ from below. Indeed, on its
positive support,
\[
 \frac{w_t(i+1)}{w_t(i)}=
 \frac{h_{i+1}/h_i}{h_{t-i}/h_{t-i-1}}\ge1
 \quad\text{if }2i+1\le t,
\]
because adjacent ratios of a positive log-concave sequence are
nonincreasing. The zero extension preserves this conclusion.
Consequently, a centered integer interval maximizes the sum of $w_t$
among sets of the same cardinality; its sum also dominates every set
of smaller cardinality.

Let $I=[L,U]\cap\mathbb Z$. In $(\sum_Ih_i)^2$, the terms on the
anti-diagonal $i+j=t$ sum over
\[
 J=I\cap(t-I),
\]
a centered interval about $t/2$.
For the product of the sums over $I-1$ and $I+1$, the corresponding set
is exactly $J-1$. The centered-interval comparison proves the sliding
inequality after summing over all $t$.

For the shrinking inequality, the neighboring intervals are
$[L-1,U+1]$ and $[L+1,U-1]$. Put $d=U-L+1$.
If $d\le2$, the latter interval is empty and the desired inequality is
immediate; if the current interval is empty, its next shrinking interval
is also empty. For $d\ge3$, the anti-diagonal cardinalities in the square
and in the neighboring product are, respectively,
\begin{align*}
 |J|&=\max\{0,\min(t-2L+1,2U-t+1,d)\},\\
 |J'|&=\max\{0,\min(t-2L+1,2U-t+1,d-2)\}.
\end{align*}
These formulas simply count pairs of integers in intervals of lengths
$d,d$, or $d+2,d-2$, whose sum is $t$. Thus $|J'|\le|J|$.
Since $J$ is centered, its $w_t$ sum is at least the sum over $J'$.
Summation over $t$ proves $V_k^2\ge V_{k-1}V_{k+1}$.
\end{proof}

\subsection*{3. Log-concavity of the lower-half first differences}
For a nonnegative integer vector $w$ (including the empty vector), set
$T=\sum_iw_i$, and write $p_w(k)=[q^k]P_w(q)$, extended by zero outside
$0\le k\le T$. Define
\begin{equation}\label{eq:gdef}
 g_w(k)=
 \begin{cases}
 p_w(k)-p_w(k-1),&0\le k\le\lfloor T/2\rfloor,\\
 0,&\text{otherwise}.
 \end{cases}
\end{equation}
The truncation is part of the definition, not an assertion about the
signs of the untruncated differences.

\begin{lemma}[Difference profile]\label{lem:g}
For every $w$, $g_w$ is nonnegative and log-concave, and is positive exactly
on an interval $\{0,1,\ldots,K_w\}$. In particular $g_w(0)=1$ and $P_w$
is symmetric and unimodal.
More explicitly, let $w=(z,m)$, $S=\sum z_i$, and $h=g_z$. Then for every
integer $k\ge0$,
\begin{equation}\label{eq:rec}
 g_{(z,m)}(k)=\sum_{i=k-m}^{\min(k,S-k)}h_i.
\end{equation}
If $\supp(h)=\{0,\ldots,r\}$, its support endpoint is
\begin{equation}\label{eq:support}
 K_{(z,m)}=\min\left\{m+r,\ S,\ \left\lfloor\frac{S+m}{2}\right\rfloor\right\}.
\end{equation}
\end{lemma}
\begin{proof}
Induct on the number of coordinates. For the empty vector, $P=1$ and
$g(0)=1$, with all other values zero. Each $P_z$ is symmetric about $S/2$,
since every factor is symmetric. Symmetry and telescoping give
\[
 p_z(k)=\sum_{i=0}^{\min(k,S-k)}h_i,
\]
where a negative upper endpoint means zero. For
$0\le k\le\lfloor(S+m)/2\rfloor$, multiplication by
$1+q+\cdots+q^m$ gives
\[
 g_{(z,m)}(k)=p_z(k)-p_z(k-m-1),
\]
which is \eqref{eq:rec}; here $k-m-1<S/2$ and is smaller than both
$k$ and $S-k$. For larger $k$, the interval in \eqref{eq:rec} is empty,
so the same identity agrees with the stipulated truncation.

By induction, $h$ is positive on $[0,r]$ and zero elsewhere, with
$r\le\lfloor S/2\rfloor$. The interval in \eqref{eq:rec} meets $[0,r]$
exactly when
\[
 k\le m+r,\qquad k\le S,\qquad 2k\le S+m.
\]
This proves nonnegativity and \eqref{eq:support}.

To prove log-concavity, put $s=\lfloor S/2\rfloor$.
For $k\le s$, \eqref{eq:rec} equals the sliding sum $U_k$ from
Lemma~\ref{lem:interval}. For $k\ge s$, it equals the shrinking sum $V_k$:
when $S=2s+1$ and $k=s$, the only extra term of $V_s$ is $h_{s+1}=0$.
Thus all log-concavity inequalities except possibly the one centered at
$s$ follow directly from that lemma. At $s$, we have
\[
 g_{(z,m)}(s-1)=U_{s-1},\quad g_{(z,m)}(s)=U_s,
 \quad g_{(z,m)}(s+1)\le U_{s+1},
\]
which proves the missing inequality as well. If $s=0$, the inequality
involving index $-1$ is automatic because that value is zero.
The induction is complete. The last assertion follows from $g_w\ge0$
and the symmetry of $P_w$.
\end{proof}

\subsection*{4. A shifted ratio that can cross a level only once}
\begin{lemma}[One-coordinate comparison]\label{lem:one}
Let $w'=w+e_j$. For $k\ge0$ on the positive support of $g_w$, the sequence
\begin{equation}\label{eq:ratio}
 \frac{g_{w'}(k+1)}{g_w(k)}
\end{equation}
is nonincreasing in $k$. Moreover, if $g_w(k)=0$ for $k\ge0$, then
$g_{w'}(k+1)=0$.
\end{lemma}
\begin{proof}
Permute coordinates so the changed coordinate is last. Write $w=(z,m)$,
$S=\sum z_i$, $h=g_z$, and $f_m=g_{(z,m)}$.
By \eqref{eq:support}, $K_{m+1}\le K_m+1$, proving the support assertion.
For positive $f_m(k)$, the denominator and numerator of \eqref{eq:ratio}
are interval sums with the same lower endpoint $L=k-m$ and respective
upper endpoints
\[
 \min(k,S-k)\quad\text{and}\quad\min(k+1,S-k-1).
\]
There are three regimes.

If $2k\le S-2$, the ratio is
\begin{equation}\label{eq:early}
 1+\frac{h_{k+1}}{\sum_{i=k-m}^k h_i}\ge1.
\end{equation}
Where $h_{k+1}>0$, the added fraction is the reciprocal of
\[
 \sum_{j=1}^{m+1}\frac{h_{k+1-j}}{h_{k+1}}.
\]
For a positive log-concave sequence, every backward ratio
$h_{t-j}/h_t$ is nondecreasing as $t$ increases. This follows by expressing
it as a product of the nondecreasing ratios $h_{a-1}/h_a$; terms below
index zero are zero until they enter the support. Thus the displayed sum
is nondecreasing, and \eqref{eq:early} is nonincreasing. Once
$h_{k+1}=0$, it stays zero as $k$ increases, and the ratio is $1$.

If $S$ is odd and $k=(S-1)/2$, the upper endpoints coincide and the
ratio is exactly $1$.

If $2k\ge S$, put $U=S-k$. The ratio is
\begin{equation}\label{eq:late}
 1-\frac{h_U}{\sum_{i=L}^U h_i}\le1.
\end{equation}
On the positive denominator support, $U\ge0$. If $h_U=0$, necessarily
$U>r$, so the ratio is initially $1$. Once $h_U>0$, the subtracted
fraction is the reciprocal of
\[
 \sum_{j=0}^{U-L}\frac{h_{U-j}}{h_U}.
\]
When $k$ increases by one, $U$ decreases by one, $L$ increases by one,
and the range of $j$ shortens by two. Each surviving backward ratio
weakly decreases by log-concavity. Hence this sum is nonincreasing, its
reciprocal is nondecreasing, and \eqref{eq:late} is nonincreasing.

The early ratios are at least $1$ and the late ratios at most $1$, with
the possible odd-parity bridge equal to $1$. This proves the comparison
across the transitions too. If the numerator reaches zero, the initial
support property prevents any later positive numerator. All divisions
above were restricted to positive denominators.
\end{proof}

\begin{lemma}[Coordinatewise comparison]\label{lem:many}
Let $v\le u$ coordinatewise and $\delta=\sum_i(u_i-v_i)$. Then, for
$s\ge0$ with $g_v(s)>0$,
\begin{equation}\label{eq:many}
 F(s)=\frac{g_u(s+\delta)}{g_v(s)}
\end{equation}
is nonincreasing in $s$. If $g_v(s)=0$ and $s\ge0$, then
$g_u(s+\delta)=0$.
\end{lemma}
\begin{proof}
Choose a path of $\delta$ unit coordinate increments
$v=w^{(0)},w^{(1)},\ldots,w^{(\delta)}=u$. Repeated support inclusion
from Lemma~\ref{lem:one} shows that a positive endpoint
$g_u(s+\delta)$ forces every $g_{w^{(a)}}(s+a)$ to be positive. On this
initial interval,
\[
 F(s)=\prod_{a=0}^{\delta-1}
 \frac{g_{w^{(a+1)}}(s+a+1)}{g_{w^{(a)}}(s+a)}.
\]
Every factor is nonnegative and nonincreasing by Lemma~\ref{lem:one},
so their product is nonincreasing. The numerator subsequently vanishes
permanently by Lemma~\ref{lem:g}, giving a zero tail. Iterating support
inclusion also proves the last assertion. If $\delta=0$, the ratio is
identically $1$ wherever defined, so the statement remains valid.
\end{proof}

\subsection*{5. The upper-half coefficients of a box ideal}
Rank the finite box
$C(u)=\prod_i\{0,1,\ldots,u_i\}$ by $|x|=\sum_i x_i$; its total rank
is $A=\sum_i u_i$. A lower ideal is a subset $J$ closed under decreasing
coordinates.

\begin{lemma}[Upper-half decrease]\label{lem:tail}
For any lower ideal $J\subseteq C(u)$ and any integer $t>A/2$,
\[
 |\{x\in J:|x|=t\}|\le |\{x\in J:|x|=t-1\}|.
\]
\end{lemma}
\begin{proof}
We first construct a saturated symmetric chain decomposition of $C(u)$.
In a rectangle $\{0,\ldots,p\}\times\{0,\ldots,m\}$, take the chain
\[
 (0,0),(1,0),\ldots,(p,0),(p,1),\ldots,(p,m).
\]
If both $p,m>0$, the remaining rectangle is
$\{0,\ldots,p-1\}\times\{1,\ldots,m\}$. Translate its second coordinate
down by one and repeat. This produces chains whose endpoint ranks sum
to $p+m$. For a product of more chains, inductively take each symmetric
chain of the earlier factors, whose endpoint ranks are $r$ and $S-r$,
and tensor it with the next chain. Apply the rectangle construction to
its length $S-2r$ and the new chain length. Shifting the ranks back by
$r$ makes every new chain symmetric about $(S+m)/2$. This proves the
required decomposition by induction.

If such a chain contains an element of rank $t>A/2$, its initial rank is
at most $A-t<t$, so this element has an immediate predecessor on the
chain. A lower ideal contains that predecessor whenever it contains the
element. Mapping each rank-$t$ element of $J$ to its predecessor is
injective, since distinct chains are disjoint. This proves the inequality.
\end{proof}

\subsection*{6. Proof of the nested-box theorem}
\begin{proof}[Proof of Theorem~\ref{thm:boxes}]
The case $u=v$ is immediate. Otherwise put
$A=\sum_i u_i$, $B=\sum_i v_i$, and $\delta=A-B>0$.
Since both products are symmetric in their respective degrees, reversal
in degree $A$ gives
\begin{equation}\label{eq:reverse}
 q^A\bigl(P_u(q^{-1})-P_v(q^{-1})\bigr)
 =P_u(q)-q^\delta P_v(q)=:H(q).
\end{equation}
Write $c_i=u_i-v_i$. The polynomial $H$ counts the lower ideal
\[
 J=C(u)\setminus
   \{x\in C(u):x_i\ge c_i\text{ for every }i\}.
\]
Indeed, the removed upper box is the translate $c+C(v)$ and has
rank-generating polynomial $q^\delta P_v(q)$. Thus $H$ is coefficientwise
nonnegative, and Lemma~\ref{lem:tail} proves that its coefficients are
nonincreasing past $A/2$.

It remains to check the first half. Let $h_t=[q^t]H(q)$ and $h_{-1}=0$.
For $0\le t\le\lfloor A/2\rfloor$, we have
$t-\delta\le B/2$; hence \eqref{eq:gdef} gives
\begin{equation}\label{eq:Hdiff}
 h_t-h_{t-1}=g_u(t)-g_v(t-\delta).
\end{equation}
For $t<\delta$, this is nonnegative. For $t\ge\delta$, set
$s=t-\delta\ge0$. Where $g_v(s)>0$, the sign of \eqref{eq:Hdiff} is
that of $F(s)-1$ from Lemma~\ref{lem:many}. Since $F$ is nonincreasing,
these signs cannot return to positive after becoming negative.
If $g_v(s)=0$, support inclusion gives $g_u(s+\delta)=0$ too, so the
difference is zero. Thus the first-half coefficients weakly increase
and then weakly decrease, with no later increase. Appending the
nonincreasing upper half from Lemma~\ref{lem:tail} proves that $H$ is
unimodal. Reversal preserves unimodality, including zero padding, and
\eqref{eq:reverse} proves the theorem.
\end{proof}
Together with Proposition~\ref{prop:rook} and Lemma~\ref{lem:g}, this proves
the exact near-full target, with no restriction on $\ell$ or the board size.

\subsection*{7. Algebraic interpretation and limits of the conclusion}
The same proof gives a characteristic-independent numerical corollary.
For a field $K$, $u\in\NN^n$, and $0\le c_i\le u_i$,
\[
 K[x_1,\ldots,x_n]/
 (x_1^{u_1+1},\ldots,x_n^{u_n+1},x_1^{c_1}\cdots x_n^{c_n})
\]
has Hilbert series
\[
 P_u(q)-q^{\sum_i c_i}P_{u-c}(q),
\]
and therefore has a unimodal Hilbert function. This follows by counting
monomials, not by asserting a Lefschetz property. For the near-full
positive-$a$ family, $c_i=1$ for every $i$.

This distinction is necessary: Migliore--Mir\'o-Roig--Nagel, Theorem 4.3,
give the non-WLP examples
$K[x_1,\ldots,x_r]/(x_1^r,\ldots,x_r^r,x_1\cdots x_r)$ for $r\ge3$.
They lie in the displayed family. Their failure of WLP is not a failure
of the numerical unimodality proved here. No assertion is made that
arbitrary-rank rook polynomials are unimodal; Lewis--Morales,
Proposition 4, disproves that different statement.

\subsection*{8. Exact computational checks and reproduction}
The accompanying \texttt{verify.py} uses only the Python standard
library and exact integers. It checks the reduction by literal placements
and tests the auxiliary lemmas. The following counts and reported passes
come from archived computation records. The complete suite was not rerun
for this edition; the archived records do not independently certify every
historical execution.
\begin{center}
\begin{tabular}{@{}p{0.66\linewidth}r@{}}
\toprule
Archived test family & Reported count\\
\midrule
Coefficient products against direct monomial enumeration & 341\\
All nonempty partitions with at most 15 cells; literal rook cancellation & 683\\
Zero polynomials among these boards & 358\\
Positive multisets of lengths 2--8, entries 1--12; near-full unimodality & 125,957\\
Nonnegative multisets of lengths 1--8, entries 0--8; $g$-log-concavity and recursion & 24,309\\
Distinct one-coordinate increments of those multisets; shifted ratio and support & 102,960\\
Deterministic broad nested-box pairs, dimensions 1--24, outer entries 0--99 & 50,000\\
Additional large one-coordinate comparisons and $g$-log-concavity tests & 10,000\\
\bottomrule
\end{tabular}
\end{center}
The broad tests also check the composed ratio of
Lemma~\ref{lem:many} directly. They include dimensions, parities, zero
coordinates, equality cases, and highly unequal inner and outer sizes.
These finite checks are diagnostics only; they are not premises of any
lemma or theorem.

To reproduce all tests, with assertions enabled, run from the package:
\begin{verbatim}
python verify.py --out-dir reproduced_results
\end{verbatim}
The deterministic generator, seed, enumeration order, and exact
coefficient computations are specified in the source.
\texttt{results/verification.json} records counts and SHA-256 digests of
canonical JSON-line record streams, including the coefficient arrays.
\texttt{results/board\_cases.json} contains every literal board result;
\texttt{results/examples.json} gives exact edge examples and the first
20 broad-test cases. Full larger streams can be regenerated exactly;
they are not required to verify the proof. The report and supplied
results correspond to the default 50,000-case run, not the optional
shortened run.

\subsection*{9. Attribution and status}
The proof was developed with AI assistance. This manuscript has not undergone
independent expert peer review or formal proof checking. No exhaustive
literature-priority claim is made for the nested-box theorem.

The product-difference reduction and monomial Hilbert-series interpretation
were present in unpublished preliminary research notes. The argument developed
here proves the interval-sum inequalities, lower-half difference-profile
log-concavity, and shifted comparisons, and combines them with the explicitly
constructed symmetric chain decomposition.

\section*{References}
\noindent
Joel Brewster Lewis and Alejandro H. Morales,
\emph{Garsia--Remmel $q$-rook numbers are not always unimodal},
arXiv:2410.19714v3. Definitions, pages 1--2; Example 3 and Proposition 4;
Remark 6, page 5.
\url{https://arxiv.org/abs/2410.19714v3}.

\noindent
Juan C. Migliore, Rosa M. Mir\'o-Roig, and Uwe Nagel,
\emph{Monomial ideals, almost complete intersections and the Weak
Lefschetz Property}, arXiv:0811.1023; \emph{Transactions of the American
Mathematical Society} \textbf{363} (2011), 229--257.
Theorem 4.3 supplies the WLP obstruction mentioned above.
\url{https://arxiv.org/abs/0811.1023}.

\noindent
Bek Chase and Filip Jonsson Kling,
\emph{The strong Lefschetz property for some monomial almost complete
intersections}, arXiv:2507.18516v1 (arXiv abstract title:
\emph{Monomial Almost Complete Intersections and the Strong Lefschetz
Property}). Relevant background: Definition 2.1, Theorem 3.1, Lemma 3.3,
and Proposition 4.2. These are not invoked to prove the unrestricted
target.
\url{https://arxiv.org/abs/2507.18516v1}.

\noindent
\emph{Unpublished preliminary research notes}: starting product-difference
reduction, Hilbert-series interpretation, and finite computations.
The reduction and Hilbert-series interpretation are derived in full above.
\end{document}
